\section{Teoría operativa del pipeline por fases}

Metodológicamente, el informe mantiene una perspectiva de conectividad funcional (no conectividad efectiva) y enfatiza la robustez frente a conducción de volumen y falsos positivos \citep{Cohen2014,Vinck2011WPLI,Cao2022EEGConnectivity}. El orden de trabajo del pipeline es intencional: ingestión de datos y controles de calidad, filtrado determinista, descomposición y limpieza por ICA, normalización por segmentos y rechazo de artefactos, CSD opcional como mejora espacial prevista, estimación de wPLI limitada por bandas y, finalmente, inferencia estadística basada en surrogates. Este orden se conserva porque cada etapa prepara un supuesto específico requerido por la siguiente; si falla un control QC, el relanzamiento debe empezar desde la fase local correspondiente.

Desde el punto de vista del desarrollo interno, el flujo se implementa como fases modulares con persistencia explícita entre pasos. Los outputs intermedios se serializan para que cualquier fase pueda relanzarse sin recomputar toda la cadena, algo esencial para depuración y análisis de sensibilidad. En la práctica, esto también permite comparar rápidamente conjuntos de parámetros (por ejemplo, umbrales de artefactos, longitud de segmento o número de surrogates) manteniendo trazabilidad completa de los artefactos generados \citep{Blinowska-Cieslak2022,Sanei2022}.
La ruta operativa puede leerse como tres bloques contiguos: armonización/adquisición y limpieza (Fases 0--2), preparación a nivel de segmento y acondicionamiento de señal (Fases 3--6), y conectividad más validación inferencial (Fases 7--8). Esta agrupación es útil al depurar porque reduce el espacio de búsqueda: los fallos antes de ICA suelen ser problemas de datos/metadatos; los fallos en el bloque medio suelen corresponder a segmentación o umbrales; y los fallos finales suelen implicar definiciones de bandas, supuestos de conectividad o configuración de surrogates. El mapa visual completo ya aparece en la \Cref{fig:generated-global-pipeline}; esta sección desarrolla su implementación fase a fase.

\subsection{Mapa del pipeline NeuroMIND y orden de ejecución}

La implementación se organiza como una cadena secuencial con persistencia local entre fases. Esto permite relanzamientos deterministas desde el punto de fallo sin recomputar las etapas previas.

\begin{tcolorbox}[
  enhanced,
  breakable,
  colback=blue!2,
  colframe=blue!55!black,
  coltitle=black,
  boxrule=0.6pt,
  arc=2mm,
  left=1.2mm,
  right=1.2mm,
  top=1mm,
  bottom=1mm,
  title=\textbf{Fases del pipeline y orden de ejecución}
]
\begin{enumerate}[leftmargin=*,itemsep=2pt,topsep=2pt]
\item \textbf{Fase 0 -- IO/BIDS, señal raw y QC inicial:} cargar EEG raw, validar metadatos, revisar canales/electrodos y calcular indicadores QC iniciales.
\item \textbf{Fase 1 -- filtrado y preprocesamiento:} aplicar filtrado determinista en orden fijo (notch 50\,Hz, bandreject 100\,Hz, high-pass 0.5\,Hz, low-pass 150\,Hz).
\item \textbf{Fase 2 -- ICA y limpieza de componentes:} descomponer en ICs, clasificar artefactos y reconstruir EEG limpio cuando proceda.
\item \textbf{Fases 3--5 -- segmentación, baseline y rechazo de artefactos:} producir épocas controladas por calidad y registrar exclusiones.
\item \textbf{Fases 6--7 -- espectro, CSD opcional, wPLI y red:} estimar potencia por banda, aplicar acondicionamiento espacial si está activo y calcular conectividad funcional.
\item \textbf{Fase 8 -- surrogates, FDR e inferencia:} construir distribuciones nulas, calcular p-valores empíricos y aplicar control de comparaciones múltiples.
\end{enumerate}
\end{tcolorbox}

El repositorio actual separa módulos fuente ejecutables de notebooks narrativos en Pluto. La Tabla~\ref{tab:pipeline-traceability} proporciona un mapa compacto de trazabilidad entre la capa de documentación publicada y los artefactos de código que ejecutan realmente cada etapa.

\begin{fullwidth}
\begin{table}[H]
\centering
\scriptsize
\caption{Mapa de trazabilidad entre fase del pipeline NeuroMIND, notebook de documentación, fuente ejecutable y principales outputs persistidos.}
\label{tab:pipeline-traceability}
\begin{adjustbox}{max width=0.98\fulltextwidth}
\begin{tabular}{@{}>{\raggedright\arraybackslash}p{0.11\fulltextwidth}>{\raggedright\arraybackslash}p{0.15\fulltextwidth}>{\raggedright\arraybackslash}p{0.29\fulltextwidth}>{\raggedright\arraybackslash}p{0.32\fulltextwidth}>{\raggedright\arraybackslash}p{0.08\fulltextwidth}@{}}
\toprule
\textbf{Fase} & \textbf{Pluto} & \textbf{Fuente ejecutable} & \textbf{Outputs persistidos} & \textbf{Estado} \\
\midrule
F0 IO/BIDS/QC & \texttt{BIDSLoader.jl}; configuración TOML & \texttt{src/io/BIDSLoader.jl}; \texttt{src/io/Config.jl}; \texttt{src/types.jl}; \texttt{src/qc/QualityControl.jl} & \bids{data/BIDS/}; \texttt{overview.csv}; \texttt{qc\_summary.csv}; \texttt{config\_snapshot.toml} & Activa \\
\addlinespace
F1 Filtrado & \texttt{Filtering.jl} & \texttt{src/preprocessing/Filtering.jl} & \texttt{filtered\_signal\_preview.png}; configuración de filtros; PSD antes/después en dashboard & Activa \\
\addlinespace
F2 ICA & núcleo, clasificación e inspección ICA & \texttt{src/ica/ICACore.jl}; \texttt{src/ica/ICAClassification.jl}; \texttt{src/ica/ICAInspection.jl} & \texttt{ica\_summary.json}; \texttt{ica\_component\_features.csv}; matrices y señales ICA & Activa \\
\addlinespace
F3--F5 Segmentos/AR & \texttt{Epochs.jl} & \texttt{src/segmentation/Epochs.jl} & \texttt{segmentation\_summary.json}; \texttt{segments\_table.csv}; \texttt{artifact\_rejection\_summary.json}; \texttt{rejected\_segments.csv} & Activa \\
\addlinespace
F6 Espectral & \texttt{PowerSpectrum.jl} & \texttt{src/spectral/PowerSpectrum.jl} & \texttt{spectral\_summary.json}; \texttt{psd\_by\_channel.csv}; \texttt{band\_power\_summary.csv}; \texttt{regional\_psd.csv} & Activa \\
\addlinespace
F7 wPLI/red & \texttt{wPLI.jl}; \texttt{CSD.jl}; \texttt{GraphMetrics.jl} & \texttt{src/connectivity/wPLI.jl}; \texttt{src/connectivity/CSD.jl}; \texttt{src/connectivity/GraphMetrics.jl} & \texttt{connectivity\_summary.json}; \texttt{connectivity\_edges.csv}; \texttt{network\_metrics.csv}; \texttt{wpli\_*.png/csv} & Activa \\
\addlinespace
F8 Inferencia & \texttt{Surrogates.jl}; \texttt{FDR.jl} & \texttt{src/statistics/Surrogates.jl}; \texttt{src/statistics/FDR.jl} & Matrices nulas surrogate, p-valores empíricos, q-valores y tablas FDR & Activa exploratoria \\
\bottomrule
\end{tabular}
\end{adjustbox}
\end{table}
\end{fullwidth}

La tabla anterior cumple ahora la función de mapeo operativo entre documentación narrativa, scripts fuente y salidas persistidas, evitando notas al pie extensas que dificulten la lectura visual del pipeline.

\subsection{Contratos operativos por fase}

Cada fase se documenta como un contrato reproducible entre teoría, código y resultados. En la práctica, esto permite revisar una ejecución sin abrir todo el código: si un resultado parece anómalo, la tabla indica qué entrada revisar, qué output persistido consultar, qué control QC aplicar y qué panel del dashboard ofrece la vista interactiva.

\begin{fullwidth}
\begin{table}[H]
\centering
\scriptsize
\caption{Contrato operativo entre fases analíticas, outputs persistidos y paneles del dashboard. Los nombres siguen el vocabulario del pipeline NeuroMIND; la numeración del dashboard es más granular que la sección metodológica.}
\label{tab:pipeline-operational-contracts}
\begin{adjustbox}{max width=0.99\fulltextwidth}
\begin{tabular}{@{}p{0.10\fulltextwidth}p{0.20\fulltextwidth}p{0.23\fulltextwidth}p{0.24\fulltextwidth}p{0.15\fulltextwidth}@{}}
\toprule
\textbf{Fase informe} & \textbf{Entrada verificable} & \textbf{Transformación} & \textbf{Outputs/QC que sostienen la fase} & \textbf{Dashboard} \\
\midrule
F0 IO/BIDS/QC & \texttt{data/BIDS/}, metadatos, electrodos & Carga, normalización sujeto/sesión/tarea, validación canal-electrodo y QC inicial & \texttt{overview.csv}, \texttt{qc\_summary.csv}, \texttt{config\_snapshot.toml} & Paneles 0--3 \\
\addlinespace
F1 Filtrado & Señal raw y configuración TOML & Notch, bandreject, high-pass y low-pass en orden determinista & Vista de señal filtrada, respuesta de filtros y PSD por etapas & Panel 4 \\
\addlinespace
F2 ICA & Señal filtrada continua & PCA whitening, FastICA simétrico, scores de artefactos y reconstrucción & \texttt{ica\_summary.json}, \texttt{ica\_component\_features.csv}, señales before/after & Panel 5 \\
\addlinespace
F3--F5 Segmentos/AR & Señal limpia ICA & Épocas, baseline, rechazo por amplitud/gradiente, cobertura por canal & Resúmenes de segmentación, tabla de épocas y segmentos rechazados & Paneles 6--7 \\
\addlinespace
F6 Espectral & Tensor de épocas válidas & Ventana, FFT, PSD y potencia por bandas/regiones & \texttt{spectral\_summary.json}, \texttt{psd\_by\_channel.csv}, \texttt{band\_power\_summary.csv}, \texttt{regional\_psd.csv} & Panel 8 \\
\addlinespace
F7 wPLI/red & Épocas por banda, CSD opcional & Señal analítica, wPLI por pares, métricas de red & \texttt{connectivity\_summary.json}, \texttt{connectivity\_edges.csv}, \texttt{network\_metrics.csv}, \texttt{wpli\_*.png/csv} & Panel 9 \\
\addlinespace
F8 Inferencia & Matrices wPLI observadas & Surrogates de fase, p empírico, Benjamini--Hochberg/FDR & Resumen surrogate, calidad nula, p/q-valores y conexiones significativas & Panel 10 \\
\addlinespace
Síntesis & Outputs por ejecución & KPIs, inventario de archivos y preparación de reporte & Log del pipeline, snapshots, categorías de archivos y copias en \texttt{report/assets/} & Paneles 11--12 \\
\bottomrule
\end{tabular}
\end{adjustbox}
\end{table}
\end{fullwidth}

\begin{controlqc}
Una fase solo debe alimentar conclusiones de la fase siguiente cuando sus outputs mínimos existen y son coherentes. Por ejemplo, wPLI debe interpretarse junto con QC, ICA, segmentación y rechazo de artefactos; la inferencia con surrogates/FDR debe interpretarse junto con las matrices observadas y los q-valores exportados.
\end{controlqc}

\subsection{Fase 0: IO/BIDS, señal raw y QC inicial (paneles 0--3)}

Los datos raw se cargan y se estandarizan en estructuras por sujeto y condición con metadatos validados. Esta fase fija etiquetas de canal, dimensiones y consistencia sujeto/condición, y genera estadísticas iniciales de calidad junto con gráficos exploratorios que sirven como línea base antes del filtrado: trazas raw, PSD promedio por canal, resúmenes de potencia por banda, gráficos de estadísticos y flags de canales sospechosos.

Para cada señal de canal $x=\{x_n\}_{n=1}^{N}$ se calculan estadísticos descriptivos que caracterizan el comportamiento basal y ayudan a detectar canales sospechosos antes del filtrado. La media resume desplazamiento DC, la desviación estándar resume dispersión, los extremos detectan clipping o saturación, RMS actúa como proxy compacto de energía/amplitud, y asimetría/curtosis describen la forma de la distribución.

\begin{parametros}
\small
\textbf{Definiciones rápidas de QC por canal.}
\[
\begin{array}{lll}
\text{Media} & \displaystyle \mu=\frac{1}{N}\sum_{n=1}^{N}x_n
& \text{nivel promedio u offset DC} \\[7pt]
\text{Desv. estándar} & \displaystyle \sigma=\sqrt{\frac{1}{N-1}\sum_{n=1}^{N}(x_n-\mu)^2}
& \text{dispersión alrededor de la media} \\[7pt]
\text{Extremos} & \displaystyle x_{\min}=\min(x),\quad x_{\max}=\max(x)
& \text{límites observados / saturación} \\[7pt]
\text{RMS} & \displaystyle \mathrm{RMS}=\sqrt{\frac{1}{N}\sum_{n=1}^{N}x_n^2}
& \text{energía o amplitud cuadrática media} \\[7pt]
\text{Asimetría} & \displaystyle \gamma_1=\frac{1}{N}\sum_{n=1}^{N}\left(\frac{x_n-\mu}{\sigma}\right)^3
& \text{sesgo de la distribución} \\[7pt]
\text{Curtosis} & \displaystyle \gamma_2=\frac{1}{N}\sum_{n=1}^{N}\left(\frac{x_n-\mu}{\sigma}\right)^4
& \text{colas pesadas / impulsividad}
\end{array}
\]
\end{parametros}
\begin{fullwidth}
\begin{figure}[H]
\centering
\resizebox{0.92\textwidth}{!}{%
\begin{tikzpicture}[x=1cm,y=1cm,>=Latex,font=\scriptsize]
  % ===== Panel 1: Mean =====
  \begin{scope}[shift={(0,2.9)}]
    \draw[black!65] (0,0) rectangle (3.9,2.1);
    \node[anchor=west,font=\scriptsize\bfseries] at (0.08,1.92) {Media};
    \draw[-Latex,black!55] (0.45,0.35)--(3.45,0.35);
    \draw[-Latex,black!55] (0.45,0.35)--(0.45,1.70);
    \node[black!65,anchor=west] at (3.00,0.42) {$t$};
    \node[black!65,rotate=90,anchor=south] at (0.18,0.52) {$\mu$V};
    \draw[black!45] (0.35,0.95)--(3.55,0.95);
    \draw[blue!70!black,thick]
      (0.35,0.75) .. controls (0.9,1.45) and (1.5,0.45) .. (2.0,1.05)
      .. controls (2.45,1.35) and (3.0,0.7) .. (3.55,1.15);
    \draw[red!75!black,dashed,thick] (0.45,1.05)--(3.55,1.05);
    \node[red!75!black,anchor=west] at (3.00,1.13) {$\mu$};
  \end{scope}

  % ===== Panel 2: Standard deviation =====
  \begin{scope}[shift={(4.2,2.9)}]
    \draw[black!65] (0,0) rectangle (3.9,2.1);
    \node[anchor=west,font=\scriptsize\bfseries] at (0.08,1.92) {Desv.\ estándar};
    \draw[-Latex,black!55] (0.45,0.35)--(3.45,0.35);
    \draw[-Latex,black!55] (0.45,0.35)--(0.45,1.70);
    \node[black!65,anchor=west] at (3.00,0.42) {$t$};
    \node[black!65,rotate=90,anchor=south] at (0.18,0.52) {$\mu$V};
    \draw[black!45] (0.4,0.45)--(0.4,1.75);
    \draw[black!45] (3.5,0.45)--(3.5,1.75);
    \draw[blue!70!black,thick]
      (0.4,0.6) .. controls (1.35,1.65) and (2.55,1.65) .. (3.5,0.6);
    \draw[red!75!black,dashed,thick] (1.95,0.45)--(1.95,1.75);
    \draw[black!70,<->] (1.95,1.25)--(2.85,1.25);
    \node[anchor=west] at (2.9,1.25) {$\sigma$};
  \end{scope}

  % ===== Panel 3: Min / Max =====
  \begin{scope}[shift={(8.4,2.9)}]
    \draw[black!65] (0,0) rectangle (3.9,2.1);
    \node[anchor=west,font=\scriptsize\bfseries] at (0.08,1.92) {Mín. / Máx.};
    \draw[-Latex,black!55] (0.45,0.35)--(3.45,0.35);
    \draw[-Latex,black!55] (0.45,0.35)--(0.45,1.70);
    \node[black!65,anchor=west] at (3.00,0.42) {$t$};
    \node[black!65,rotate=90,anchor=south] at (0.18,0.52) {$\mu$V};
    \draw[black!45] (0.35,0.95)--(3.55,0.95);
    \draw[blue!70!black,thick]
      (0.35,0.95)--(0.9,1.55)--(1.35,0.55)--(1.9,1.1)--(2.45,0.35)--(3.0,1.35)--(3.55,0.8);
    \fill[red!80!black] (2.45,0.35) circle (1.4pt);
    \fill[green!60!black] (0.9,1.55) circle (1.4pt);
    \node[red!80!black,anchor=west] at (2.52,0.38) {$x_{\min}$};
    \node[green!50!black,anchor=west] at (0.97,1.58) {$x_{\max}$};
  \end{scope}

  % ===== Panel 4: RMS =====
  \begin{scope}[shift={(0,0)}]
    \draw[black!65] (0,0) rectangle (3.9,2.1);
    \node[anchor=west,font=\scriptsize\bfseries] at (0.08,1.92) {RMS};
    \draw[-Latex,black!55] (0.45,0.35)--(3.45,0.35);
    \draw[-Latex,black!55] (0.45,0.35)--(0.45,1.70);
    \node[black!65,anchor=west] at (3.00,0.42) {$t$};
    \node[black!65,rotate=90,anchor=south] at (0.18,0.52) {$\mu$V};
    \draw[black!45] (0.35,0.95)--(3.55,0.95);
    \draw[blue!70!black,thick]
      (0.35,1.35).. controls (0.9,0.5) and (1.35,1.55) .. (1.9,0.8)
      .. controls (2.35,0.45) and (2.95,1.45) .. (3.55,0.7);
    \draw[orange!85!black,dashed,thick] (0.35,0.35)--(3.55,0.35);
    \node[orange!85!black,anchor=west] at (2.45,0.45) {RMS};
  \end{scope}

  % ===== Panel 5: Skewness =====
  \begin{scope}[shift={(4.2,0)}]
    \draw[black!65] (0,0) rectangle (3.9,2.1);
    \node[anchor=west,font=\scriptsize\bfseries] at (0.08,1.92) {Asimetría};
    \draw[-Latex,black!55] (0.45,0.35)--(3.45,0.35);
    \draw[-Latex,black!55] (0.45,0.35)--(0.45,1.70);
    \node[black!65,anchor=west] at (3.00,0.42) {$t$};
    \node[black!65,rotate=90,anchor=south] at (0.18,0.52) {$\mu$V};
    \draw[black!45] (0.4,0.45)--(0.4,1.75);
    \draw[black!45] (3.5,0.45)--(3.5,1.75);
    \draw[blue!70!black,thick]
      (0.4,0.6) .. controls (1.0,1.6) and (1.9,1.45) .. (2.5,0.95)
      .. controls (3.0,0.65) and (3.25,0.58) .. (3.5,0.55);
    \node[anchor=west] at (2.42,1.55) {cola derecha};
  \end{scope}

  % ===== Panel 6: Kurtosis =====
  \begin{scope}[shift={(8.4,0)}]
    \draw[black!65] (0,0) rectangle (3.9,2.1);
    \node[anchor=west,font=\scriptsize\bfseries] at (0.08,1.92) {Curtosis};
    \draw[-Latex,black!55] (0.45,0.35)--(3.45,0.35);
    \draw[-Latex,black!55] (0.45,0.35)--(0.45,1.70);
    \node[black!65,anchor=west] at (3.00,0.42) {$t$};
    \node[black!65,rotate=90,anchor=south] at (0.18,0.52) {$\mu$V};
    \draw[black!45] (0.4,0.45)--(0.4,1.75);
    \draw[black!45] (3.5,0.45)--(3.5,1.75);
    \draw[blue!70!black,thick]
      (0.4,0.5) .. controls (1.35,1.9) and (2.55,1.9) .. (3.5,0.5);
    \draw[black!45,dashed] (0.4,0.75) .. controls (1.35,1.45) and (2.55,1.45) .. (3.5,0.75);
    \node[anchor=west] at (2.35,1.58) {pico alto};
  \end{scope}
\end{tikzpicture}
}
\caption{Intuición visual de los estadísticos básicos usados en la Fase 0: media, dispersión, extremos, energía RMS, asimetría y curtosis.}
\end{figure}
\end{fullwidth}
En términos operativos de QC, canales con varianza o desviación estándar anormalmente bajas pueden indicar mal contacto o registros planos, mientras que curtosis muy alta o valores mínimo/máximo extremos pueden sugerir artefactos transitorios.

\begin{fullwidth}
\begin{figure}[H]
\centering
\resizebox{0.92\fulltextwidth}{!}{\NSRawQCFlow}
\caption{Flujo de control de calidad para EEG raw reconstruido en TikZ: las métricas descriptivas de Fase 0 no se interpretan de forma aislada, sino como una cadena de detección, revisión visual y decisión antes de propagar datos al filtrado.}
\label{fig:raw-qc-threshold-flow}
\end{figure}
\end{fullwidth}

\subsection{Fase 1: filtrado y preprocesamiento (panel 4)}

El preprocesamiento toma la salida de Fase 0 (IO): señales EEG organizadas como diccionario o matriz indexada por canal, con metadatos temporales/de canal y estadísticas opcionales de inspección inicial. El objetivo es retirar ruido de red, interferencias relacionadas con hardware y deriva lenta, además de limitar el ancho de banda antes del análisis de conectividad por bandas.

\begin{fullwidth}
\begin{figure}[H]
\centering
\resizebox{0.92\fulltextwidth}{!}{\NSFilterCascade}
\caption{Cascada determinista de filtrado. Cada etapa produce una salida auditable para comparar respuesta del filtro, PSD antes/después y efecto temporal sobre la señal antes de propagar datos a ICA.}
\label{fig:generated-filter-cascade}
\end{figure}
\end{fullwidth}

La implementación en \path{src/Preprocessing/filtering.jl} carga datos serializados desde \path{data/Preprocessing/IO/}, aplica una cascada de filtros digitales y, en cada paso, almacena el resultado bajo \path{data/Preprocessing/filtering/}, grafica la respuesta del filtro (magnitud y fase) y compara PSD promedio antes y después del filtrado.

Todos los filtros son Butterworth. Las etapas high-pass y low-pass usan \texttt{filtfilt} (filtrado de fase cero), por lo que el orden efectivo se duplica y se preserva la alineación temporal. La cascada se aplica en este orden:

\begin{enumerate}
\item \textbf{Notch (50\,Hz):} etapa bandstop para interferencia de red eléctrica en Europa, orden de diseño 4 y ancho de 1\,Hz alrededor de 50\,Hz.
\item \textbf{Bandreject (100\,Hz):} etapa bandstop estrecha (orden 4, ancho 1\,Hz) para atenuar posible interferencia de hardware cerca de 100\,Hz sin recortar de forma amplia frecuencias altas.
\item \textbf{High-pass (0.5\,Hz):} eliminación de deriva lenta preservando actividad baja; orden de diseño 4 (efectivo 8 con \texttt{filtfilt}).
\item \textbf{Low-pass (150\,Hz):} limitación superior de ancho de banda; orden de diseño 4 (efectivo 8 con \texttt{filtfilt}).
\end{enumerate}

Las salidas por etapa se guardan para preservar trazabilidad y permitir relanzamientos aislados de ICA o análisis espectrales posteriores.

\begin{parametros}
Orden operativo fijo: notch 50\,Hz $\rightarrow$ bandreject 100\,Hz $\rightarrow$ high-pass 0.5\,Hz $\rightarrow$ low-pass 150\,Hz. El orden se mantiene estable para que las comparaciones entre ejecuciones dependan de parámetros explícitos y no de cambios implícitos en la secuencia.
\end{parametros}

\subsection{Fase 2: ICA y limpieza de componentes (panel 5)}

Después del preprocesamiento (Fase 1), ICA se aplica para separar actividad neural de fuentes artefactuales (ocular, muscular, ruido de línea y transitorios abruptos). La descomposición sigue el modelo lineal estándar
\[
\mathbf{X}=\mathbf{A}\mathbf{S},\qquad \mathbf{S}=\mathbf{W}\mathbf{X},
\]
con $\mathbf{X}\in\mathbb{R}^{C\times N}$ (canales $\times$ muestras), $\mathbf{A}$ la matriz de mezcla y $\mathbf{W}$ la matriz de demixing estimada por FastICA simétrico.

\begin{fullwidth}
\begin{figure}[H]
\centering
\resizebox{0.9\fulltextwidth}{!}{\NSICAFlow}
\caption{Flujo ICA generado desde cero: señales mezcladas, separación en componentes, identificación de familias artefactuales y reconstrucción del EEG limpio con componentes retenidos.}
\label{fig:generated-ica-flow}
\end{figure}
\end{fullwidth}
\sidenote{%
\footnotesize
\textbf{ICA implementation}\\
Fase 2a: \texttt{src/ICA/ICA.jl}\\
Fase 2b: \texttt{src/ICA/ICA\_cleaning.jl}\\
Entradas desde \texttt{data/Preprocessing/filtering/}\\
Salidas a \texttt{data/Processing/ICA/}
}

En la Fase 2a, el objetivo es estimar componentes independientes y sus mapas espaciales a partir del EEG filtrado, preservando la dimensionalidad completa de canales ($k=C$). La entrada es la matriz filtrada $\mathbf{X}\in\mathbb{R}^{C\times N}$ procedente de la Fase 1.

\begin{decisionmetodologica}
\small
Los apuntes internos de ICA refuerzan una decisión importante del pipeline: ICA se trata como separación ciega de fuentes para limpieza y auditoría, no como localización cortical. El procedimiento intenta separar fuentes estadísticamente independientes mezcladas linealmente en los sensores; después, la decisión de retirar componentes se basa en evidencia convergente de topografía, espectro, estadísticos y trazas temporales.
\end{decisionmetodologica}

La primera operación es el centrado, que elimina offsets medios por canal y produce $\mathbf{X}_c$. Después se aplica blanqueamiento PCA para que los datos transformados tengan covarianza identidad y el algoritmo ICA pueda centrarse en estructura de orden superior:
\[
\mathbf{Z}=\mathbf{V}_{\text{whit}}\mathbf{X}_c,\qquad \mathbf{V}_{\text{whit}}=\mathbf{\Lambda}^{-1/2}\mathbf{E}^{\top}.
\]
Aquí, $\mathbf{E}$ contiene los autovectores de la matriz de covarianza y $\mathbf{\Lambda}$ los autovalores correspondientes. Tras el blanqueamiento, FastICA simétrico itera con $g(u)=\tanh(au)$ y decorrelación simétrica hasta convergencia.

Operativamente, el blanqueamiento también actúa como \emph{sphering}: si la covarianza del EEG centrado es $\mathbf{C}=\mathbf{E}\mathbf{\Lambda}\mathbf{E}^{\top}$, la transformación $\mathbf{\Lambda}^{-1/2}\mathbf{E}^{\top}$ reescala los ejes principales para que la covarianza de $\mathbf{Z}$ sea aproximadamente identidad. Desde ese espacio esférico, ICA estima una matriz de rotación/desmezcla que maximiza independencia aproximada mediante no gaussianidad.

La salida de esta fase es la matriz de componentes $\mathbf{S}$, la matriz total de demixing $\mathbf{W}_{\text{total}}=\mathbf{W}\mathbf{V}_{\text{whit}}$ y la matriz de mezcla $\mathbf{A}=\mathbf{W}_{\text{total}}^{-1}$, que proyecta la actividad de los ICs de vuelta al espacio de canales. En algunos entornos esta matriz de mezcla se denomina \emph{unmixing/back-projection} según la convención de orientación; en el informe se mantiene la notación $\mathbf{A}$ para reconstrucción y $\mathbf{W}$ para separación.
\sidenote{%
\scriptsize
\textbf{FastICA outputs}\\
$\mathbf{S}$: independent components\\
$\mathbf{W}_{\text{total}}=\mathbf{W}\mathbf{V}_{\text{whit}}$\\
$\mathbf{A}=\mathbf{W}_{\text{total}}^{-1}$
}

\begin{table}[H]
\centering
\small
\caption{Lectura técnica de la Fase 2a a partir de la formulación ICA usada en NeuroMIND.}
\label{tab:ica-theory-operational}
\begin{tabular}{@{}p{0.20\textwidth}p{0.34\textwidth}p{0.36\textwidth}@{}}
\toprule
\textbf{Elemento} & \textbf{Papel matemático} & \textbf{Implicación práctica en EEG} \\
\midrule
Centrado & $\mathbf{X}_c=\mathbf{X}-\mathbb{E}[\mathbf{X}]$ por canal. & Retira offsets DC antes de calcular covarianza y evita que la media domine la descomposición. \\
Whitening / sphering & PCA/SVD de la covarianza y reescalado por autovalores. & Reduce redundancia lineal y deja al algoritmo ICA la búsqueda de estructura no gaussiana. \\
FastICA & Iteración de punto fijo con decorrelación simétrica. & Busca componentes independientes aproximados; la convergencia depende de tolerancia, semilla y número máximo de iteraciones. \\
Back-projection & Reconstrucción $\mathbf{X}_{\text{clean}}=\mathbf{A}\mathbf{S}_{\text{clean}}$. & Devuelve la señal limpia al espacio de canales para fases posteriores de segmentación, espectro y conectividad. \\
\bottomrule
\end{tabular}
\end{table}

\begin{controlqc}
\small
\textbf{Supuestos y cautelas de ICA.} La interpretación de ICs presupone mezcla lineal instantánea, independencia estadística aproximada y fuentes suficientemente no gaussianas. La escala, el signo y el orden de los componentes son indeterminados; por eso un IC no se interpreta por su índice numérico, sino por su patrón conjunto de topografía, PSD, curso temporal y puntuaciones automáticas.
\end{controlqc}

\subsubsection*{Lectura de topoplots ICA}

Un mapa topográfico de ICA, o \emph{topoplot}, no es un heatmap convencional de filas y columnas. Es una representación de la cabeza vista desde arriba: la nariz aparece en la parte superior, las orejas quedan a izquierda y derecha, y los puntos marcan la posición aproximada de los electrodos. El color no representa ``actividad cerebral'' directa, sino el peso con el que un componente independiente se proyecta sobre cada sensor a través de la matriz de mezcla $\mathbf{A}$ \citep{Cohen2014,Sanei2022}. Por tanto, la pregunta práctica no es ``dónde se activa el cerebro?'', sino ``la geometría de este componente parece compatible con una fuente cerebral o con un artefacto?''.

La escala de color expresa polaridad y magnitud relativa. Un lóbulo rojo intenso y un lóbulo azul intenso, separados de forma suave y espacialmente coherente, sugieren un patrón dipolar fuerte; regiones cercanas al blanco indican poca contribución de ese IC sobre esos electrodos. El signo rojo/azul no tiene valor fisiológico absoluto, porque ICA puede invertir el signo de un componente sin cambiar el modelo. Lo relevante es la forma: simetría, suavidad, localización frontal, distribución periférica, concentración posterior o presencia de focos muy estrechos.

\begin{table}[H]
\centering
\small
\caption{Guía operativa para interpretar topoplots ICA en NeuroMIND. La decisión final debe combinar mapa espacial, espectro, curso temporal y scores automáticos.}
\label{tab:ica-topoplot-reading-guide}
\begin{tabular}{@{}p{0.24\textwidth}p{0.36\textwidth}p{0.30\textwidth}@{}}
\toprule
\textbf{Familia de componente} & \textbf{Firma topográfica esperada} & \textbf{Lectura práctica} \\
\midrule
Parpadeo / VEOG & Máximos frontales bilaterales o mancha muy intensa en Fp1/Fp2, normalmente con energía lenta. & Candidato claro a artefacto ocular vertical; revisar traza temporal y potencia $<4$\,Hz. \\
Movimiento ocular horizontal / HEOG & Dos focos frontolaterales de polaridad opuesta cerca de F7/F8 o regiones temporales anteriores. & Candidato a movimiento izquierda--derecha; suele rechazarse si la traza confirma sacadas. \\
EMG muscular & Focos pequeños, muy intensos y periféricos, a menudo temporales o cervical/occipitales. & Candidato a músculo; confirmar con energía alta y dispersa por encima de 20--30\,Hz. \\
Ruido de línea & Topografía a veces poco fisiológica, con patrón espacial débil o irregular; decisión dominada por pico espectral en 50/60\,Hz. & No basta el mapa: el marcador fuerte es el espectro del IC y su score de línea. \\
Salto/transitorio & Foco localizado o patrón inestable, frecuentemente asociado a un sensor o borde del montaje. & Candidato a pop de electrodo, movimiento o discontinuidad; revisar curso temporal. \\
Alfa posterior & Distribución suave, relativamente simétrica y posterior, centrada en O1/O2/Oz/Pz. & Compatible con actividad cerebral de reposo EC si aparece pico alrededor de 8--12\,Hz. \\
Mu / sensorimotor & Dipolo lateral o central cerca de C3/C4, a veces diagonal. & Compatible con ritmo sensorimotor; contrastar con banda alfa/beta baja. \\
Frontal/medio cognitivo & Distribución suave alrededor de Fz/Cz, sin focos periféricos abruptos. & Posible fuente neural o mezcla residual; requiere evidencia espectral y temporal. \\
\bottomrule
\end{tabular}
\end{table}

\begin{decisionmetodologica}
\small
En NeuroMIND la clasificación ICA se formula como una revisión convergente. Los topoplots aportan la geometría espacial; los scores aportan evidencia cuantitativa; la PSD distingue familias como línea o EMG; y la traza temporal confirma parpadeos, sacadas o saltos. Herramientas externas de etiquetado automático, como ICLabel en otros ecosistemas, siguen esta misma lógica general: combinan forma espacial, frecuencia y dinámica temporal. En la implementación actual de NeuroMIND la decisión queda documentada por los ficheros \texttt{ica\_component\_features.csv}, \texttt{ica\_components.csv} y las figuras exportadas por componente.
\end{decisionmetodologica}

\begin{algorithm}[H]
\caption{Pseudocódigo operativo de FastICA simétrico (Fase 2a)}
\begin{algorithmic}
\Require $\mathbf{X}\in\mathbb{R}^{C\times N}$, $k=C$, $\texttt{max\_iter}$, $\tau$, $a$, seed
\State Center channels: $\mathbf{X}_c \gets \mathbf{X}-\mathrm{mean}(\mathbf{X},\text{dims}=2)$
\State PCA whitening: $\mathbf{Z}\gets \mathbf{\Lambda}^{-1/2}\mathbf{E}^{\top}\mathbf{X}_c$
  \State Inicializar $\mathbf{W}\in\mathbb{R}^{k\times k}$ aleatoriamente y decorrelar simétricamente
\For{$t=1$ to $\texttt{max\_iter}$}
  \State $\mathbf{Y}\gets \mathbf{W}\mathbf{Z}$
  \State $\mathbf{G}\gets \tanh(a\mathbf{Y})$, $\mathbf{G}'\gets a\,(1-\tanh^2(a\mathbf{Y}))$
  \State $\mathbf{D}\gets \operatorname{diag}(\mathrm{mean}(\mathbf{G}',\text{dims}=2))$
  \State $\mathbf{W}_{\text{new}}\gets \frac{1}{N}\mathbf{G}\mathbf{Z}^{\top}-\mathbf{D}\mathbf{W}$
  \State Symmetric decorrelation of $\mathbf{W}_{\text{new}}$
  \State $\Delta\gets \max_i\left|1-\left|\left(\mathbf{W}_{\text{new}}\mathbf{W}^{\top}\right)_{ii}\right|\right|$
  \If{$\Delta<\tau$}
    \State $\mathbf{W}\gets \mathbf{W}_{\text{new}}$; \textbf{break}
  \EndIf
  \State $\mathbf{W}\gets \mathbf{W}_{\text{new}}$
\EndFor
\State $\mathbf{S}\gets \mathbf{W}\mathbf{Z}$, $\mathbf{W}_{\text{total}}\gets \mathbf{W}\mathbf{V}_{\text{whit}}$, $\mathbf{A}\gets \mathbf{W}_{\text{total}}^{-1}$
\Ensure $\mathbf{S}$, $\mathbf{W}_{\text{total}}$, $\mathbf{A}$
\end{algorithmic}
\end{algorithm}

En la Fase 2b, cada IC se evalúa para distinguir actividad neural de actividad artefactual antes de la reconstrucción. La etapa toma como entrada $(\mathbf{S},\mathbf{W}_{\text{total}},\mathbf{A})$, metadatos de canal y, opcionalmente, coordenadas de electrodos para inspección topográfica. La regla operativa sigue el protocolo RS-MIND: el número de componentes se mantiene igual al número de canales seleccionados cuando procede, y la detección automática propone candidatos que deben poder revisarse de forma semiautomática/manual mediante topografía, PSD, curso temporal y scores.

Las características espaciales, espectrales y estadísticas se estandarizan (z-score) y se combinan en cuatro puntuaciones orientadas a artefactos: ocular, muscular, línea y salto. El índice global de decisión se define como
\[
\begin{aligned}
\mathrm{artifact\_score}_i
&=
\max\!\big(
\mathrm{ocular\_score}_i,\,
\mathrm{muscle\_score}_i,\\
&\qquad
\mathrm{line\_score}_i,\,
\mathrm{jump\_score}_i
\big).
\end{aligned}
\]
Así, cada componente se juzga por su familia de características más sospechosa. Si un solo tipo de artefacto es alto, el IC se marca como sospechoso aunque las demás puntuaciones sean bajas. Los ICs por encima del umbral se etiquetan como artefacto y se eliminan poniendo a cero las filas correspondientes en $\mathbf{S}$.

\begin{controlqc}
\small
\textbf{Lectura operativa de scores ICA.}
\par\medskip
\begin{minipage}[t]{0.62\textwidth}
\scriptsize
\begin{tabular}{@{}>{\ttfamily\bfseries}p{0.22\textwidth}p{0.36\textwidth}@{}}
ocular\_score & Parpadeo/EOG: predominio frontal, energía lenta y, si existe canal EOG, correlación con EOG. \\
muscle\_score & EMG: potencia alta en bandas rápidas y trazas de alta frecuencia. \\
line\_score & Contaminación de red: picos en 50/60 Hz y armónicos dentro del IC. \\
jump\_score & Transitorios abruptos: saltos, picos o clipping parcial, normalmente no fisiológicos. \\
\end{tabular}
\end{minipage}\hfill
\begin{minipage}[t]{0.30\textwidth}
\centering
\resizebox{\linewidth}{!}{%
\begin{tikzpicture}[x=1cm,y=1cm,>=Latex,font=\scriptsize,baseline={(0,0)}]
  \draw[black!55,->] (0,0) -- (0,2.4);
  \draw[black!55,->] (0,0) -- (5.7,0);
  \draw[orange!85!black,thick,dashed] (0,1.5) -- (5.45,1.5);
  \node[orange!85!black,anchor=west] at (3.35,1.72) {umbral 1.5};
  \foreach \x/\h/\lab/\col in {
    0.75/0.65/ocular/teal,
    2.05/1.95/muscular/orange,
    3.35/0.55/línea/purple,
    4.65/0.85/salto/blue}
  {
    \fill[\col!70] (\x,0) rectangle +(0.55,\h);
    \node[rotate=35,anchor=east] at (\x+0.35,-0.12) {\lab};
  }
  \draw[red!75!black,thick,->] (2.32,2.15) -- (2.32,1.98);
  \node[red!75!black,anchor=west] at (2.55,2.12) {$\max$};
  \node[anchor=west] at (3.15,1.1) {IC sospechoso};
  \draw[red!75!black,thick,->] (5.25,0.85) -- (4.95,0.85);
\end{tikzpicture}}
\end{minipage}
\end{controlqc}

El EEG limpio se reconstruye entonces mediante
\[
\mathbf{X}_{\text{clean}}=\mathbf{A}\mathbf{S}_{\text{clean}}.
\]
Esta ecuación restaura la señal en el espacio de canales usando solo los componentes retenidos. Las salidas finales incluyen el EEG limpio indexado por canal y metadatos completos de auditoría (lista de ICs malos, tabla de características, etiquetas y comparaciones opcionales) para preservar reproducibilidad y trazabilidad.

\begin{parametros}
\textbf{Umbral práctico.} La regla típica es \texttt{artifact\_score > 1.5}. Un valor más bajo aumenta la sensibilidad y rechaza más ICs; un valor más alto es más conservador y retiene componentes dudosos para revisión manual.
\end{parametros}


\subsection{Fases 3--5: segmentación, baseline y rechazo de artefactos (paneles 6--7)}

Tras la limpieza ICA (Fase 2b), el EEG continuo se segmenta en épocas de longitud fija para construir una representación estable para análisis espectral y de conectividad. La implementación actual usa segmentos de \SI{1}{s} sin solapamiento, alineados con el protocolo RS-MIND para reposo. La segmentación convierte cada serie temporal de canal en un tensor 3D

\begin{fullwidth}
\begin{figure}[H]
\centering
\resizebox{0.9\fulltextwidth}{!}{\NSEpochBaseline}
\caption{Segmentación y baseline como flujo de control: señal continua, épocas de duración fija, corrección de baseline, rechazo de artefactos y tensor limpio con métricas QC asociadas.}
\label{fig:generated-epochs-baseline-flow}
\end{figure}
\end{fullwidth}

\[
\mathbf{E}\in\mathbb{R}^{C\times L\times K},
\]
donde $C$ es el número de canales, $L$ el número de muestras por época y $K$ el número de épocas válidas.
\begin{figure}[H]
\centering
\includegraphics[width=0.54\textwidth]{methodology/segmentacion_epoch_ventana_2s.pdf}
\caption{Concepto de época: ventana temporal corta y de duración fija extraída de la señal continua.}
\label{fig:margin-epoch-window}
\end{figure}
\begin{figure}[H]
\centering
\resizebox{0.62\columnwidth}{!}{%
\begin{tikzpicture}[>=Latex]
  % Front face
  \draw[rounded corners=1.5pt, draw=black!75, fill=black!3] (0,0) rectangle (3.2,2.1);
  % Back face offset
  \draw[draw=black!55, fill=black!8] (0.9,0.7) rectangle (4.1,2.8);
  % Connect corners
  \draw[black!55] (0,0)--(0.9,0.7);
  \draw[black!55] (3.2,0)--(4.1,0.7);
  \draw[black!55] (0,2.1)--(0.9,2.8);
  \draw[black!55] (3.2,2.1)--(4.1,2.8);
  % Inner hint lines
  \draw[black!25] (1.6,0)--(1.6,2.1);
  \draw[black!25] (0,1.05)--(3.2,1.05);
  \node at (2.15,1.55) {$\mathbf{E}$};
  % Axis arrows
  \draw[-Latex,thick] (-0.55,0)--(-0.55,2.1);
  \node[anchor=east] at (-0.62,1.05) {$C$};
  \draw[-Latex,thick] (0,-0.45)--(3.2,-0.45);
  \node[below] at (1.6,-0.45) {$L$};
  \draw[-Latex,thick] (3.2,2.1)--(4.4,3.05);
  \node[anchor=west] at (4.45,3.05) {$K$};
\end{tikzpicture}
}
\caption{Representación tensorial de épocas ($C \times L \times K$).}
\end{figure}

La longitud de época se define a partir de la frecuencia de muestreo y la duración objetivo:
\[
L=\left\lfloor F_s\cdot T_{\mathrm{seg}}\right\rfloor.
\]
Por ejemplo, con $T_{\mathrm{seg}}=1\,\mathrm{s}$ y $F_s=500\,\mathrm{Hz}$, cada época contiene $L=500$ muestras. La distancia inicio-a-inicio entre épocas consecutivas es
\[
\mathrm{step}=L-\mathrm{overlap}.
\]
Con solapamiento cero, $\mathrm{step}=L$, de modo que las épocas son consecutivas y no comparten muestras.

Dada una señal con $N_{\mathrm{total}}$ muestras, el número de épocas completas es
\[
K=\left\lfloor\frac{N_{\mathrm{total}}-L}{\mathrm{step}}\right\rfloor+1.
\]
Las muestras finales que no completan una época se descartan. Esto evita ventanas de longitud variable y mantiene una forma tensorial consistente entre canales y sujetos.

El valor por defecto de 1 segundo es un compromiso práctico: proporciona una resolución espectral
\[
\Delta f\approx\frac{F_s}{L}\approx 1\,\mathrm{Hz},
\]
adecuada para conectividad basada en bandas, manteniendo una hipótesis razonable de cuasi-estacionariedad dentro de cada época.

Operativamente, \path{src/Processing/segmentation.jl} carga el EEG limpio desde \path{data/Processing/ICA/}, realiza el corte, calcula estadísticos descriptivos por época (media, desviación estándar, mínimo, máximo, RMS), produce gráficos de ejemplo (épocas superpuestas y apiladas) y guarda tensores y metadatos bajo \path{data/Processing/segmentation/}. Estas salidas se corrigen por baseline y pasan al rechazo de artefactos basado en amplitud, con una segunda corrección de baseline opcional después del rechazo.

\paragraph{Corrección de baseline (antes de artefactos).}
Tras la segmentación, cada época puede seguir conteniendo un offset DC o una deriva local lenta distinta entre épocas. La primera corrección de baseline elimina este offset restando, en cada época, el valor medio de una ventana inicial predefinida. Si la ventana de baseline en segundos es $[t_{\mathrm{start}},t_{\mathrm{end}}]$, sus índices de muestra son
\[
n_{\mathrm{start}}=\lfloor t_{\mathrm{start}}F_s\rfloor+1,\qquad
n_{\mathrm{end}}=\lfloor t_{\mathrm{end}}F_s\rfloor.
\]
Para el canal $c$ y la época $k$, con $N_{\mathrm{bl}}=n_{\mathrm{end}}-n_{\mathrm{start}}+1$, la media de baseline es
\[
\mu_{c,k}=\frac{1}{N_{\mathrm{bl}}}\sum_{n=n_{\mathrm{start}}}^{n_{\mathrm{end}}}E_{c,n,k},
\]
y la época corregida es
\[
E^{\mathrm{corr}}_{c,:,k}=E_{c,:,k}-\mu_{c,k}.
\]
Esto conserva la forma y la varianza dentro de la época, pero centra la ventana de baseline alrededor de cero, reduciendo sesgo en métricas de amplitud y espectrales. En este proyecto, la ventana por defecto son los primeros 100\,ms ($t_{\mathrm{start}}=0$, $t_{\mathrm{end}}=0.1$\,s), que a $F_s=500$\,Hz corresponden a las muestras 1--50.

\begin{fullwidth}
\begin{figure}[H]
\centering
\includegraphics[width=0.88\linewidth]{methodology/baseline_correction_proceso.pdf}
\caption{Proceso visual de corrección de baseline: eliminación progresiva de offset y deriva hasta obtener una señal estable para el análisis.}
\label{fig:baseline-correction-neuromind}
\end{figure}
\end{fullwidth}

\begin{figure}[H]
\centering
\resizebox{0.86\columnwidth}{!}{%
\begin{tikzpicture}[x=1cm,y=1cm,>=Latex,font=\scriptsize]
  % Axes
  \draw[->, gray!70] (-2.8,-1.8) -- (4.5,-1.8)
    node[right] {tiempo relativo al evento (ms)};
  \draw[->, gray!70] (-2.8,-1.6) -- (-2.8,1.4)
    node[above, rotate=90, anchor=south] {EEG amplitude ($\mu$V)};

  % Time ticks
  \draw[gray!65] (-2,-1.86)--(-2,-1.74);
  \draw[gray!65] (0,-1.86)--(0,-1.74);
  \node[below] at (-2,-1.8) {$-200$};
  \node[below] at (0,-1.8) {$0$};

  % Event marker
  \draw[dashed, gray] (0,-1.8)--(0,1.15);
  \node[anchor=west] at (0.08,1.08) {inicio de evento / estímulo};

  % Baseline window region
  \fill[blue!8] (-2, -1.25) rectangle (0, 1.05);
  \draw[blue!35, dashed] (-2,-1.25) rectangle (0,1.05);
  \node[anchor=west] at (-1.9,0.95) {ventana de baseline};

  % Shared waveform shape and corrected offset
  \draw[blue!75!black, line width=1.05pt, domain=-2.4:4.0, samples=120]
    plot (\x,{0.28 + 0.23*sin(95*\x) + 0.10*cos(185*\x)});
  \draw[teal!75!black, line width=1.05pt, domain=-2.4:4.0, samples=120]
    plot (\x,{-0.62 + 0.23*sin(95*\x) + 0.10*cos(185*\x)});

  % Baseline mean in pre-event window
  \draw[dashed, gray!70] (-2,0.28)--(0,0.28);
  \node[anchor=east] at (-2.05,0.28) {$\mu_{\mathrm{baseline}}$};
  \node[anchor=west, gray!70] at (-1.95,0.12) {media de la ventana};

  % Gentle horizontal guides for levels
  \draw[gray!30] (-2.6,0.0)--(4.2,0.0);
  \draw[gray!30] (-2.6,-0.90)--(4.2,-0.90);

  % Correction arrow
  \draw[-Latex, very thick, black!70] (4.25,0.28)--(4.25,-0.62);
  \node[anchor=west] at (4.33,-0.17) {restar $\mu_{\mathrm{baseline}}$};

  % Curve labels
  \node[anchor=west] at (2.35,0.78) {antes};
  \node[anchor=west] at (2.35,-0.14) {después};

  % Formula in clear area
  \node[anchor=west] at (0.55,-1.35)
    {$x_{\mathrm{corrected}}(t)=x(t)-\mu_{\mathrm{baseline}}$};
\end{tikzpicture}
}
\caption{Efecto conceptual de la corrección de baseline en una época: desplazamiento vertical constante usando la media de la ventana de baseline.}
\end{figure}

\paragraph{Por qué corregir baseline antes y después del rechazo de artefactos?}
La primera corrección de baseline (antes del rechazo de artefactos) estandariza offsets para que las reglas por umbral de amplitud sean comparables entre épocas y canales. Sin este paso, épocas con niveles DC distintos podrían rechazarse de más o de menos por razones puramente relacionadas con offset.

Después del rechazo, el conjunto de épocas retenidas ya no es idéntico al conjunto previo, y la eliminación de épocas puede desplazar la tendencia central del dataset. Una segunda corrección opcional recentra el tensor limpio final usado para FFT/CSD/wPLI, mejorando la consistencia entre sujetos y relanzamientos. En síntesis: la primera baseline ayuda a tomar decisiones justas de artefactos; la segunda armoniza la entrada final del análisis.

\paragraph{Rechazo de artefactos basado en amplitud.}
Incluso después de ICA y baseline, algunas épocas pueden incluir ráfagas residuales (por ejemplo, picos EMG, pops de electrodo o saturaciones breves). El rechazo en \texttt{src/Processing/artifact\_rejection.jl} aplica una regla conservadora de amplitud y elimina épocas completas que violan los límites configurados.

\begin{figure}[H]
\centering
\includegraphics[width=0.54\textwidth]{methodology/rechazo_artefactos_revision.pdf}
\caption{Revisión visual de épocas contaminadas por parpadeo o movimiento antes de excluirlas del tensor final.}
\label{fig:margin-epoch-artifact-review}
\end{figure}

Sea $\mathbf{E}^{\mathrm{bl}}\in\mathbb{R}^{C\times L\times K}$ el tensor de épocas corregidas por baseline y sea
\[
\mathcal{C}_{\mathrm{used}}=\{1,\dots,\min(N_{\mathrm{used}},C)\}
\]
el conjunto de canales inspeccionados para aplicar umbrales (por defecto $N_{\mathrm{used}}=30$). Para cada época $k$, se definen el mínimo y máximo global sobre canales inspeccionados y muestras:
\[
m_k=\min_{c\in\mathcal{C}_{\mathrm{used}},\,1\le n\le L}E^{\mathrm{bl}}_{c,n,k},
\qquad
M_k=\max_{c\in\mathcal{C}_{\mathrm{used}},\,1\le n\le L}E^{\mathrm{bl}}_{c,n,k}.
\]
La época $k$ se marca como mala si
\[
m_k<A_{\min}\quad\text{or}\quad M_k>A_{\max},
\]
con límites conservadores por defecto $A_{\min}=-70\,\mu\mathrm{V}$ y $A_{\max}=+70\,\mu\mathrm{V}$, coherentes con el protocolo operativo RS-MIND para esta fase.
\sidenote{%
\footnotesize
\textbf{Lógica de umbral (una época)}\\[1pt]
\centering
\begin{tikzpicture}[x=0.95cm,y=0.55cm,>=Latex]
  \draw[black!35] (0,0)--(4.8,0);
  \draw[red!70!black, dashed, line width=0.7pt] (0,0.85)--(4.8,0.85);
  \draw[red!70!black, dashed, line width=0.7pt] (0,-0.85)--(4.8,-0.85);
  \node[anchor=west, scale=0.68] at (4.85,0.85) {$A_{\max}$};
  \node[anchor=west, scale=0.68] at (4.85,-0.85) {$A_{\min}$};
  \draw[blue!75!black, line width=1.0pt]
    (0.1,0.10) .. controls (1.1,0.45) and (2.2,-0.20) .. (3.0,0.25)
              .. controls (3.5,0.55) and (3.8,1.10) .. (4.25,0.65);
  \node[scale=0.70] at (2.0,1.18) {violación};
\end{tikzpicture}\\
\textit{rechazar si alguna muestra cruza los límites}
}

Las épocas retenidas forman el tensor filtrado $\mathbf{E}^{\mathrm{AR}}$ concatenando solo índices válidos. Una métrica de calidad útil es la tasa de rechazo
\[
\mathrm{rate}=\frac{K-K_{\mathrm{AR}}}{K},
\]
donde $K_{\mathrm{AR}}$ es el número de épocas retenidas. Para diagnóstico, el conjunto de canales violadores
\[
\mathcal{V}_k=
\left\{
c\in\mathcal{C}_{\mathrm{used}}:\,
\min_n E^{\mathrm{bl}}_{c,n,k}<A_{\min}
\;\text{or}\;
\max_n E^{\mathrm{bl}}_{c,n,k}>A_{\max}
\right\}
\]
se guarda junto con estadísticos mínimo/máximo por época.
\sidenote{%
\footnotesize
\textbf{Valores de procesamiento por defecto}\\
Segment length: \texttt{1 s}\\
Overlap: \texttt{0}\\
AR thresholds: \texttt{$\pm$70 $\mu$V}\\
Used channels for AR: \texttt{30}\\
Baseline window: \texttt{0--0.1 s}\\
Segunda baseline: opcional
}

\subsection{Fases 6--7: análisis espectral, CSD opcional, wPLI y red funcional (paneles 8--9)}

\begin{fullwidth}
\begin{figure}[H]
\centering
\resizebox{0.92\fulltextwidth}{!}{\NSCSDWPLI}
\caption{CSD + wPLI como decisión metodológica central prevista: CSD mejora especificidad espacial a nivel de sensores cuando está activo, y wPLI resume sincronización de fase no nula reduciendo el peso de acoplamientos de desfase cero.}
\label{fig:generated-csd-wpli}
\end{figure}
\end{fullwidth}

\begin{decisionmetodologica}
La fase de conectividad no busca reconstruir fuentes corticales ni inferir direccionalidad. El diseño metodológico contempla CSD previo para reducir contaminación espacial; sin embargo, la ejecución preliminar documentada en resultados usa \texttt{use\_csd=false}. Por tanto, esas matrices se leen como conectividad funcional sensor--sensor para validación operativa, no como salida final CSD+wPLI.
\end{decisionmetodologica}

\subsubsection*{7.1 Transformación Current Source Density (CSD)}
La conducción de volumen propaga corrientes a través de los tejidos de la cabeza, de modo que cada electrodo registra una mezcla ponderada de múltiples fuentes. Esto infla el acoplamiento de desfase cero entre canales cercanos y reduce la especificidad espacial. Para mejorar la interpretación antes de estimar conectividad, el pipeline aplica una transformación Current Source Density (CSD), es decir, una aproximación de Laplaciano superficial de la segunda derivada espacial del potencial en cuero cabelludo \citep{Cohen2014}.

Operativamente, \texttt{src/Connectivity/CSD.jl} carga segmentos corregidos por segunda baseline y coordenadas de electrodos (BIDS TSV), comprueba consistencia canal-electrodo y aplica un Laplaciano de splines esféricos estilo Perrin (equivalente BVA). Sea
\[
\mathbf{E}\in\mathbb{R}^{C\times L\times K}
\]
el tensor segmentado y $\mathbf{L}\in\mathbb{R}^{C\times C}$ el operador CSD. Para cada segmento $s$,
\[
\mathbf{E}_{\mathrm{CSD}}[:,:,s]=\mathbf{L}\,\mathbf{E}[:,:,s].
\]
Esto produce una representación efectivamente libre de referencia (la componente constante se elimina por construcción), más adecuada para conectividad basada en fase.

Las coordenadas de electrodos se normalizan primero a la esfera unidad:
\[
r_c=\sqrt{x_c^2+y_c^2+z_c^2},\qquad
\mathbf{R}_c=\left(\frac{x_c}{r_c},\frac{y_c}{r_c},\frac{z_c}{r_c}\right).
\]
\begin{figure}[H]
\centering
\resizebox{0.78\columnwidth}{!}{%
\begin{tikzpicture}[font=\scriptsize]
  % Head outline
  \draw[thick] (0,0) circle (2.6);
  % Nose
  \draw[thick] (-0.35,2.58) -- (0,2.95) -- (0.35,2.58);
  % Ears
  \draw[thick] (-2.6,0.65) .. controls (-2.95,0.25) and (-2.95,-0.25) .. (-2.6,-0.65);
  \draw[thick] (2.6,0.65) .. controls (2.95,0.25) and (2.95,-0.25) .. (2.6,-0.65);
  % Axes
  \draw[-Latex,black!70] (-3.2,-3.0) -- (-1.8,-3.0) node[below] {$x$};
  \draw[-Latex,black!70] (-3.2,-3.0) -- (-3.2,-1.6) node[left] {$y$};

  % Electrode positions (2D approximation in normalized head coordinates)
  \foreach \name/\x/\y in {
    Fp1/-0.95/2.05, Fp2/0.95/2.05,
    F7/-2.00/1.20, F3/-0.95/1.25, Fz/0/1.35, F4/0.95/1.25, F8/2.00/1.20,
    T7/-2.25/0.15, C3/-1.05/0.20, Cz/0/0.25, C4/1.05/0.20, T8/2.25/0.15,
    P7/-1.90/-1.00, P3/-0.95/-0.85, Pz/0/-0.90, P4/0.95/-0.85, P8/1.90/-1.00,
    O1/-0.75/-2.00, Oz/0/-2.10, O2/0.75/-2.00,
    TP9/-2.25/-1.55, TP10/2.25/-1.55
  }{
    \filldraw[blue!65!black] (\x,\y) circle (1.2pt);
    \node[black] at (\x,\y+0.19) {\name};
  }
\end{tikzpicture}
}
\caption{Disposición aproximada 2D del cuero cabelludo (vista x-y) con electrodos etiquetados usados para operadores espaciales como CSD.}
\end{figure}

Usando productos escalares $\cos\gamma_{ij}=\mathbf{R}_i\cdot\mathbf{R}_j$, los kernels spline se construyen mediante expansiones truncadas de Legendre (hasta \texttt{degree}):
\[
G_{ij}=\sum_{k=1}^{\texttt{degree}} c_g(k)\,P_k(\cos\gamma_{ij}),
\qquad
H_{ij}=\sum_{k=1}^{\texttt{degree}} c_h(k)\,P_k(\cos\gamma_{ij}).
\]
Tras la regularización $G_\lambda=G+\lambda I$, la restricción libre de referencia se impone con
\[
 C
 =
 G_\lambda^{-1}
 -
 \frac{(G_\lambda^{-1}\mathbf{1})(\mathbf{1}^{\top}G_\lambda^{-1})}
      {\mathbf{1}^{\top}G_\lambda^{-1}\mathbf{1}},
\]
y el operador CSD final es
\[
\mathbf{L}=H\,C.
\]
\sidenote{%
\scriptsize
\textbf{CSD defaults (BVA-aligned)}\\
$m$ (spline order): \texttt{4}\\
\texttt{degree}: \texttt{10}\\
$\lambda$: $10^{-5}$\\
Script: \texttt{src/Connectivity/CSD.jl}
}

\begin{algorithm}[H]
\caption{CSD mediante Laplaciano de splines esféricos de Perrin (como en \texttt{CSD.jl})}
\begin{algorithmic}
\Require $\mathbf{E}\in\mathbb{R}^{C\times L\times K}$, BIDS electrode TSV $(x_c,y_c,z_c)$, $m$, \texttt{degree}, $\lambda$
  \State Cargar filas de electrodos tipo EEG y ajustar orden a la lista de canales
\If{any channel lacks coordinates}
  \State abortar con error de consistencia
\EndIf
\For{$c=1$ to $C$}
  \State Normalizar posición a esfera unidad: $\mathbf{R}_c\gets (x_c,y_c,z_c)/\sqrt{x_c^2+y_c^2+z_c^2}$
\EndFor
\State Precalcular coeficientes de Perrin $c_g(k), c_h(k)$ para $k=1,\dots,\texttt{degree}$
\For{$i=1$ to $C$}
  \For{$j=i+1$ to $C$}
    \State $\cos\gamma \gets \mathrm{clamp}(\mathbf{R}_i\!\cdot\!\mathbf{R}_j,\,-1,1)$
    \State Construir $G_{ij}$ y $H_{ij}$ mediante sumas de Legendre
  \EndFor
\EndFor
\State $G_\lambda\gets G+\lambda I$; enforce reference-free constraint to compute $C$
\State $\mathbf{L}\gets H\,C$
\For{$s=1$ to $K$}
  \State $\mathbf{E}_{\mathrm{CSD}}[:,:,s]\gets \mathbf{L}\,\mathbf{E}[:,:,s]$
\EndFor
\Ensure $\mathbf{E}_{\mathrm{CSD}}$ y artefactos CSD exportados (matrices/figuras/tablas)
\end{algorithmic}
\end{algorithm}

\subsubsection*{7.2 Descomposición espectral y señales limitadas por banda}
Después de la segunda corrección opcional de baseline (Fase 6), el EEG segmentado se usa para estimar espectros en el dominio frecuencial y derivar descriptores de potencia por banda para conectividad y control de calidad. La rutina espectral (\texttt{src/Spectral/FFT.jl}) sigue una secuencia compatible con BVA sobre el tensor segmentado y aplica FFT por canal/segmento antes de promediar posteriormente por banda y época:
\[
\mathbf{E}^{\mathrm{bl2}}\in\mathbb{R}^{C\times L\times K}.
\]

\begin{figure}[H]
\centering
\resizebox{0.88\textwidth}{!}{\NSFFTFlow}
\caption{Flujo conceptual usado en la Sección 7.2 para descomposición espectral basada en FFT y estimación de potencia por banda.}
\end{figure}

Para cada canal $c$ y segmento $k$, se elimina una estimación DC por segmento antes de aplicar ventana,
\[
\mu_{c,k}=\frac{1}{L}\sum_{n=1}^{L}E^{\mathrm{bl2}}_{c,n,k},
\qquad
E^{\mathrm{dc}}_{c,:,k}=E^{\mathrm{bl2}}_{c,:,k}-\mu_{c,k},
\]
y después se aplica una ventana tipo Hamming/Tukey $w[n]$:
\[
E^{\mathrm{win}}_{c,:,k}=E^{\mathrm{dc}}_{c,:,k}\odot w.
\]
El zero-padding hasta $N_{\mathrm{fft}}=512$ produce el tensor listo para FFT $\mathbf{E}^{\mathrm{pad}}$, con resolución frecuencial
\[
\Delta f=\frac{F_s}{N_{\mathrm{fft}}}\approx 0.976\ \mathrm{Hz}\quad (F_s=500\ \mathrm{Hz}).
\]

La FFT real se calcula por canal/segmento:
\[
X_{c,:,k}=\mathrm{rFFT}\!\left(E^{\mathrm{pad}}_{c,:,k}\right),
\]
con $n_{\mathrm{bins}}=N_{\mathrm{fft}}/2+1$. La componente DC se reintroduce en el bin 0 (convención BVA) y la potencia por bin se obtiene como
\[
P_{c,\ell,k}=|X_{c,\ell,k}|^2,
\]
seguida de corrección por varianza de ventana (división por $\overline{w^2}$), plegado de espectro completo (duplicando bins interiores) y normalización por $N_{\mathrm{fft}}^2$, de modo que la potencia espectral final se expresa en $\mu V^2$.

La potencia por banda se calcula promediando bins cuya frecuencia central cae dentro de cada intervalo $[f_{\min},f_{\max})$ (Delta, Theta, Alpha, Beta baja/media/alta, Gamma), primero por segmento y luego entre segmentos.
\sidenote{%
\scriptsize
\textbf{FFT defaults}\\
$N_{\mathrm{fft}}=512$\\
$\Delta f\approx 0.976$ Hz\\
Window: Hamming/Tukey\\
Unidad de salida: $\mu V^2$
}

\begin{algorithm}[H]
\caption{Análisis espectral mediante FFT (como en \texttt{FFT.jl})}
\begin{algorithmic}
\Require $\mathbf{E}^{\mathrm{bl2}}\in\mathbb{R}^{C\times L\times K}$, $F_s$, $N_{\mathrm{fft}}=512$, window parameters
\For{$c=1$ to $C$}
  \For{$k=1$ to $K$}
    \State $\mu_{c,k}\gets \mathrm{mean}(\mathbf{E}^{\mathrm{bl2}}[c,:,k])$
    \State $\mathbf{E}^{\mathrm{dc}}[c,:,k]\gets \mathbf{E}^{\mathrm{bl2}}[c,:,k]-\mu_{c,k}$
    \State Build taper $w[n]$ and apply $\mathbf{E}^{\mathrm{win}}[c,:,k]\gets \mathbf{E}^{\mathrm{dc}}[c,:,k]\odot w$
    \State Zero-pad to length $N_{\mathrm{fft}}$: $\mathbf{E}^{\mathrm{pad}}[c,1\!:\!L,k]\gets \mathbf{E}^{\mathrm{win}}[c,:,k]$
    \State $X_{c,:,k}\gets \mathrm{rFFT}(\mathbf{E}^{\mathrm{pad}}[c,:,k])$
    \State Reintroducir DC en bin 0 y calcular potencia $P_{c,:,k}\gets |X_{c,:,k}|^2$
    \State Aplicar corrección de varianza, plegado espectral y normalización por $N_{\mathrm{fft}}^2$
  \EndFor
\EndFor
\State Para cada banda $b$, promediar $P$ sobre bins con $f_\ell\in[f_{\min}^{(b)},f_{\max}^{(b)})$
\Ensure Potencia por bin, espectros promediados por segmento y salidas de potencia por banda
\end{algorithmic}
\end{algorithm}

\subsubsection*{7.3 Métricas de conectividad wPLI}
La conectividad funcional basada en fase se estima con weighted Phase Lag Index (wPLI), que enfatiza relaciones de fase no nula consistentes y reduce el peso de contribuciones cercanas a desfase cero, típicamente asociadas a conducción de volumen. Para cada banda de frecuencia, las señales de entrada de la fase de conectividad (CSD si está activo; sensores corregidos si \texttt{use\_csd=false}) se filtran en banda y se convierten en representaciones analíticas mediante transformada de Hilbert; después, wPLI se agrega sobre todas las muestras y segmentos.

\begin{fullwidth}
\begin{figure}[H]
\centering
\resizebox{0.92\fulltextwidth}{!}{\NSWPLIIntuition}
\caption{Intuición de wPLI inspirada en las figuras vectoriales de \texttt{wpli\_vector\_figures}: el estimador penaliza contribuciones compatibles con desfase cero y resume la consistencia del signo de la parte imaginaria del espectro cruzado en una matriz de conectividad por banda.}
\label{fig:generated-wpli-intuition}
\end{figure}
\end{fullwidth}

Sea
\[
S_{xy}(t)=z_i(t)\,\overline{z_j(t)}
\]
el término cruzado entre señales analíticas en los canales $i$ y $j$, y definimos
\[
\mathrm{Im}_{ij}(t)=\operatorname{Im}\!\left(S_{xy}(t)\right)=\operatorname{Im}\!\left(z_i(t)\overline{z_j(t)}\right).
\]
Entonces wPLI se define como
\[
\mathrm{wPLI}_{ij}
=
\frac{\left|\mathbb{E}\left[\mathrm{Im}_{ij}(t)\right]\right|}
{\mathbb{E}\left[\left|\mathrm{Im}_{ij}(t)\right|\right]+\epsilon}.
\]
Solo se usa la componente imaginaria porque el acoplamiento de desfase cero tiende a dominar la parte real bajo conducción de volumen o referencia común. Frente a PLI simple (que usa solo signos), wPLI pondera cada muestra por $|\mathrm{Im}_{ij}(t)|$, reduciendo sensibilidad a diferencias de fase débiles y ruidosas y mejorando robustez en muestras finitas \citep{Vinck2011WPLI}.

La matriz de conectividad $\mathbf{W}\in\mathbb{R}^{C\times C}$ es simétrica, tiene diagonal cero y valores en $[0,1]$. Los resultados se exportan por banda/condición como matrices, tablas de aristas, heatmaps y logs.
\sidenote{%
\scriptsize
\textbf{wPLI defaults}\\
Band-pass order: \texttt{8}\\
Filtering: \texttt{filtfilt} (zero-phase)\\
Aggregation: all $L\times K$ samples\\
$\epsilon$: \texttt{eps()}
}

\begin{algorithm}[H]
\caption{Cálculo de wPLI por banda (como en \texttt{wPLI.jl})}
\begin{algorithmic}
\Require $\mathbf{E}_{\mathrm{conn}}\in\mathbb{R}^{C\times L\times K}$, band $(f_1,f_2)$, $F_s$
\For{$c=1$ to $C$}
  \For{$s=1$ to $K$}
    \State $x\gets \mathbf{E}_{\mathrm{conn}}[c,:,s]$ \Comment{CSD si activo; sensores si \texttt{use\_csd=false}}
    \State $x_f\gets \mathrm{bandpass\_filtfilt}(x,F_s,f_1,f_2,\mathrm{order}=8)$
    \State $Z_c[:,s]\gets \mathrm{analytic\_signal}(x_f)$ \Comment{Hilbert vía FFT}
  \EndFor
\EndFor
\For{$i=1$ to $C$}
  \For{$j=i+1$ to $C$}
    \State $\mathrm{Im}_{ij}\gets \operatorname{Im}(Z_i\odot \overline{Z_j})$
    \State $\mathrm{num}\gets |\sum \mathrm{Im}_{ij}|$;\; $\mathrm{den}\gets \sum |\mathrm{Im}_{ij}|+\epsilon$
    \State $W[i,j]\gets \mathrm{num}/\mathrm{den}$;\; $W[j,i]\gets W[i,j]$
  \EndFor
\EndFor
\Ensure $\mathbf{W}$ (simétrica, diagonal 0), matriz CSV, aristas CSV y heatmap PNG
\end{algorithmic}
\end{algorithm}

\subsection{Fase 8: surrogates, FDR e inferencia (panel 10)}

Después de obtener matrices empíricas de conectividad, la fase final construye distribuciones nulas basadas en surrogates y realiza inferencia estadística con control de comparaciones múltiples. En la ejecución revisada de NeuroMIND, esta fase ya exporta matrices de p-valores, q-valores, máscaras significativas y resúmenes nulos por banda. La especificación siguiente documenta la lógica algorítmica usada y conserva el marco de auditoría necesario para relanzar la inferencia con una configuración final congelada.

\begin{fullwidth}
\begin{figure}[H]
\centering
\resizebox{0.92\fulltextwidth}{!}{\NSSurrogatesFDR}
\caption{Lógica inferencial de surrogates y FDR: la aleatorización de fase conserva potencia por canal, destruye relaciones de fase entre canales y permite comparar cada arista real contra una distribución nula empírica antes de aplicar FDR.}
\label{fig:generated-surrogates-fdr}
\end{figure}
\end{fullwidth}

\begin{pendiente}
Aunque la fase ya produce salidas inferenciales para la ejecución revisada, todavía debe tratarse como exploratoria hasta fijar configuración, CSD, número de surrogates, semilla y convención de exportación para ejecuciones grupales reproducibles.
\end{pendiente}

\subsubsection{8.1 Modelo de datos surrogate}
Para evaluar significación, se generan señales EEG surrogate de forma independiente por canal mediante aleatorización de fase en el dominio de Fourier. Para cada canal, si $X(f)$ es la transformada de Fourier de la señal original, los espectros surrogate se definen como
\[
X^{(s)}(f)=|X(f)|\,e^{\mathrm{i}\phi^{(s)}(f)},
\qquad
\phi^{(s)}(f)\sim\mathcal{U}(0,2\pi),
\]
forzando simetría conjugada para que la FFT inversa produzca señales reales.

Esta construcción preserva la densidad espectral de potencia por canal,
\[
\left|X^{(s)}(f)\right|^2=\left|X(f)\right|^2,
\]
mientras destruye relaciones de fase entre canales. Bajo esta nula, la conectividad observada se atribuye a alineamiento de fase por azar y no a acoplamiento genuino.

Operativamente, cada realización surrogate repite la misma cadena de la Fase 7 (filtrado por banda, extracción de señal analítica y cálculo wPLI), produciendo matrices surrogate de conectividad por banda y condición.

\subsubsection{8.2 Marco de inferencia estadística}
Sea $\mathrm{wPLI}^{\mathrm{real}}_{ij}$ el valor empírico de la arista y $\mathrm{wPLI}^{(s)}_{ij}$ las realizaciones surrogate para la arista $(i,j)$. El p-valor empírico es
\[
p_{ij}
=
\frac{1+\#\{\mathrm{wPLI}^{(s)}_{ij}\ge\mathrm{wPLI}^{\mathrm{real}}_{ij}\}}
{1+N_s},
\]
donde $N_s$ es el número de surrogates. Esta forma con suma de uno evita p-valores cero en muestras surrogate finitas y proporciona un soporte nulo empírico discreto.

Los p-valores se ajustan después por multiplicidad entre aristas usando False Discovery Rate (FDR, Benjamini--Hochberg; $\alpha=0.05$ por defecto). La interpretación final se basa, por tanto, en aristas que sobreviven tanto a la comparación contra la nula empírica como al control de multiplicidad.
\sidenote{%
\footnotesize
\textbf{Inference knobs}\\
\texttt{N\_s}: number of surrogates\\
Correction: FDR (default)\\
Trade-off: runtime vs p-value precision
}
